\documentclass[11pt]{article}

%================ Packages ================
\usepackage[utf8]{inputenc}
\usepackage[T1]{fontenc}
\usepackage{amsmath}
\usepackage[a4paper,margin=1in]{geometry}
\usepackage{xcolor}
\usepackage{mdframed}
\usepackage{booktabs}
\usepackage{tabularx}
\usepackage{enumitem}
\usepackage{titlesec}
\usepackage{lastpage}
\usepackage{fancyhdr}
\usepackage{natbib}
\usepackage[colorlinks=true,linkcolor=accent,urlcolor=accent,citecolor=accent]{hyperref}

%================ Colours ================
\definecolor{accent}{HTML}{1F5FA8}
\definecolor{revgray}{HTML}{F4F5F7}
\definecolor{revrule}{HTML}{B9C2CE}
\definecolor{respgreen}{HTML}{2E7D46}

%================ Header / footer ================
\pagestyle{fancy}
\fancyhf{}
\renewcommand{\headrulewidth}{0.4pt}
\lhead{\small\textcolor{accent}{\textbf{Minomaly} -- Response to Reviewers}}
\rhead{\small KAIS}
\cfoot{\small\thepage\ / \pageref{LastPage}}

%================ Section styling ================
\titleformat{\section}{\large\bfseries\color{accent}}{\thesection.}{0.5em}{}
\titleformat{\subsection}{\normalsize\bfseries}{\thesubsection}{0.5em}{}

%================ Reviewer-comment box ================
\newmdenv[
  backgroundcolor=revgray, linecolor=revrule, linewidth=0.6pt,
  topline=false, bottomline=false, rightline=false, leftline=true,
  innerleftmargin=10pt, innerrightmargin=10pt,
  innertopmargin=6pt, innerbottommargin=6pt,
  skipabove=8pt, skipbelow=4pt
]{reviewerbox}
\newcommand{\comment}[1]{\begin{reviewerbox}\itshape #1\end{reviewerbox}}
\newcommand{\response}{\par\noindent\textbf{\textcolor{respgreen}{Response.}}\ }
\setlist[itemize]{leftmargin=1.4em,itemsep=2pt,topsep=3pt}

%================ Title ================
\title{\vspace{-1.2cm}\color{accent}\Large\textbf{Response to the Editor and Reviewers}}
\author{}\date{}

\begin{document}
\maketitle
\vspace{-1.6cm}
\begin{center}
\begin{tabular}{rl}
\textbf{Manuscript:} & \emph{Minomaly: A Subgraph Mining Framework for Unsupervised}\\
 & \emph{and Interpretable Anomaly Detection in Graphs}\\[2pt]
\textbf{Journal:} & Knowledge and Information Systems (KAIS)\\
\textbf{Prior decision:} & Reject\\
\end{tabular}
\end{center}
\noindent\rule{\textwidth}{0.8pt}

%==================================================================
\section{Letter to the Editor}
\noindent Dear Prof.\ Wu,\\[4pt]

We thank you and the reviewers for the assessment. We were encouraged that the reviewers found the framework \emph{``technically coherent''} and \emph{``complete''} (R1), described the interpretability as \emph{``a genuine contribution''} with \emph{``compelling''} extracted-pattern tables (R3), and called the pattern-mining/GNN bridge \emph{``a principled formulation''} with \emph{``clean theoretical grounding''} (R3).

We respectfully appeal the decision and have prepared a revision that addresses every substantive point. The handling editor raised four concerns; we summarize our response here and detail it below.

\begin{enumerate}[leftmargin=1.6em,itemsep=3pt]
  \item \textbf{Limited novelty over SPMiner.} This was the most important and most frequently raised point (Editor; R1.W1; R3.W1). The submitted version relied on SPMiner's pretrained model. \textbf{The revision removes this dependence}: we introduce \textbf{DSAN}, an order-embedding subgraph matcher that we train ourselves (GLASS-style labeling, softplus order-embedding loss). Minomaly now reuses only the order-embedding \emph{principle}; the model, the labeling, and the loss are our own, and DSAN measurably improves detection (Table on SPMiner~vs.~DSAN).
  \item \textbf{Missing recent / relevant baselines.} We added the dense-subgraph fraud-detection line that R3 identified --- \textbf{Fraudar} (KDD~2016), \textbf{Spade} (VLDB~2023), \textbf{RUSH} (VLDB~2024) --- to the related work, positioned them precisely against our setting, and validate Minomaly on a real-world fraud graph (see point~3).
  \item \textbf{Only artificially injected anomalies.} We added experiments with a \textbf{different, non-clique injection} (DICE-structural) so the evaluation no longer rests on a single anomaly type that could match our inductive bias. We also added an experiment on the \textbf{Amazon fraud detection graph} \citep{dou2020enhancing}, a real-world dataset with organic fraud labels, where Minomaly discovers the fraud cliques as rare structures and the structural interpretation of the returned patterns identifies them with 94.1\% AUC (details in R3.W4).
  \item \textbf{Only three datasets.} We added a \textbf{fourth dataset (CiteSeer)} with seven injection ratios (0.5\%--8\%) and a \textbf{fifth dataset (Amazon fraud)} with organic labels, for a total of five evaluation settings.
\end{enumerate}

We believe these revisions are concrete and complete, and we hope the manuscript may be reconsidered. A point-by-point response follows.

\vspace{6pt}\noindent Sincerely,\\ The Authors

%==================================================================
\section{Response to the Handling Editor}
\comment{The paper offers limited novelty over SPMiner [28]. The reviewers point out relevant recent related baselines which are not included in the related work or experimental comparison. There are only experiments on data with artificially injected anomalies. Also, only 3 data sets are used.}
\response We have addressed all four points: (i) the SPMiner dependence is removed and replaced by our own trained model, DSAN (see R1.W1/R3.W1); (ii) the dense-subgraph fraud baselines Fraudar, Spade, and RUSH are now discussed in the related work, and we validate Minomaly on the real-world Amazon fraud graph (see R3.W4); (iii) we added a non-clique injection (DICE-structural) and an organic-label fraud experiment (see R3.W3/R3.W4); and (iv) we added CiteSeer (seven injection ratios) and Amazon fraud (organic labels) for a total of five evaluation settings. Details follow under each reviewer.

%==================================================================
\section{Response to Reviewer 1}

\subsection*{R1.W1 -- Novelty positioning relative to SPMiner}
\comment{The paper does not explain clearly enough what is inherited from earlier order-embedding / subgraph-mining work, what is adapted for anomaly detection, and what is new.}
\response We agree, and we made a substantive change rather than a rewording. The submitted version relied on SPMiner's pretrained model; the revision \textbf{introduces DSAN}, an order-embedding matcher we train ourselves from scratch, removing the SPMiner dependence. We now state explicitly what is inherited (the order-embedding \emph{principle}), what is adapted (order embeddings repurposed from frequent-motif mining to estimating each node's neighborhood-containment frequency), and what is new (DSAN with a GLASS-style labeling trick and a smooth softplus loss; the anti-monotone frequency measure; and the greedy anomalous-structure search). The phrase ``we rely on SPMiner for subgraph mining'' has been removed.

\subsection*{R1.W2 -- Related work organization and gap}
\comment{The related work mixes topics and does not build a clear path to the problem; reorganize into subsections ending with a gap statement.}
\response The Related Work is reorganized into \emph{Graph Anomaly Detection} $\rightarrow$ \emph{Interpretable Anomaly Detection} $\rightarrow$ \emph{Pattern Mining and Order Embedding}, ending with an explicit gap statement that leads into Minomaly. The order-embedding line most related to our work now appears last, immediately before our positioning.

\subsection*{R1.W3 -- Overuse of ``state-of-the-art''}
\comment{The repeated ``state-of-the-art'' is not well supported; baselines are from 2016--2022.}
\response We softened every occurrence to ``competitive/established baselines'' and now describe the baselines explicitly as deep-learning GAD methods from 2016--2022.

\subsection*{R1.W4 -- Clarity / over-long sentences, and minor issues}
\comment{Several sentences are too long (e.g.\ the SPMiner sentence). Minor: ``Graph a Anomaly Detection''; reference [48] ``Erd6s'' should be ``Erd\H{o}s''.}
\response We rewrote the flagged SPMiner sentence and other long sentences into shorter statements, and fixed both editorial issues (``Graph Anomaly Detection''; ``Erd\H{o}s'').

%==================================================================
\section{Response to Reviewer 3}

\subsection*{R3.W1 -- Core mechanism borrowed from SPMiner}
\comment{The OES construction --- GNN architecture, hinge loss, order-embedding framework --- is taken from SPMiner; ``We rely on SPMiner''; the pretrained model is SPMiner's artifact. The contribution is thin.}
\response This is the same concern as R1.W1 and the editor's, and it is the one we changed most. We no longer use SPMiner's model. We train our own matcher, \textbf{DSAN}, which differs from SPMiner in two concrete ways: a \textbf{GLASS-style labeling trick} (anchor + inside-subgraph indicators) that provably increases the structural expressiveness of the GNN, and a \textbf{softplus order-embedding loss} that replaces the hinge and gives sharper subgraph/non-subgraph separation. Table~\ref{tab:dsan_vs_spminer_resp} reports an ablation with the search procedure held fixed. DSAN improves detection on Cora (AUC $95.3\rightarrow97.1$, AP $74.6\rightarrow94.4$) and Amazon (AUC $95.4\rightarrow97.6$, Recall $92.9\rightarrow97.1$), and matches SPMiner on Flickr, while running 2--9$\times$ faster. Together with the anomaly-detection formulation (the anti-monotone frequency measure) and the greedy search, the contribution is no longer a thin application of SPMiner.

\begin{table}[h]
\footnotesize
\centering
\caption{Effect of replacing the pretrained SPMiner matcher with DSAN (ours), search procedure fixed. DSAN improves Cora and Amazon and matches Flickr.}
\label{tab:dsan_vs_spminer_resp}
\begin{tabular}{llccccc}
\toprule
\textbf{Dataset} & \textbf{Matcher} & \textbf{AUC} & \textbf{F1} & \textbf{AP} & \textbf{Recall} & \textbf{Precision} \\
\midrule
\multirow{2}{*}{Cora}   & SPMiner          & 95.3 & 87.1 & 74.6 & 91.4 & 83.1 \\
                         & \textbf{DSAN}    & \textbf{97.1} & \textbf{97.1} & \textbf{94.4} & \textbf{94.3} & \textbf{100.0} \\
\addlinespace
\multirow{2}{*}{Amazon}  & SPMiner          & 95.4 & 69.6 & 57.7 & 92.9 & 55.7 \\
                         & \textbf{DSAN}    & \textbf{97.6} & 68.5 & 57.3 & \textbf{97.1} & 52.9 \\
\addlinespace
\multirow{2}{*}{Flickr}  & SPMiner          & 90.5 & 87.3 & 77.7 & 81.2 & 94.3 \\
                         & \textbf{DSAN}    & \textbf{91.0} & 87.3 & 77.7 & 82.1 & 93.3 \\
\bottomrule
\end{tabular}
\end{table}

\subsection*{R3.W2 -- Missing complexity analysis}
\comment{There is no formal time/space complexity analysis for the framework or Algorithm 1, despite the scalability claim.}
\response We added a complexity analysis. OES construction embeds $k$ sampled neighborhoods at $O(k \cdot c_{\text{GNN}})$ cost, where $c_{\text{GNN}}$ is the (linear-in-edges) GraphSAGE forward pass; frequency estimation for a query is a single batched comparison against the $k$ references. The greedy search costs $O(|S| \cdot \texttt{max\_steps} \cdot \texttt{max\_cands} \cdot c_{\text{GNN}})$ for $|S|$ starting nodes, and the heuristic starting-node selection keeps $|S|$ far below $|V|$. We also relate these terms to the measured wall-clock times.

\subsection*{R3.W3 -- Only three datasets, all with synthetic clique injection}
\comment{All datasets use synthetic injection; the method's bias may align with injected cliques. No organic anomalies, heterogeneous graphs, or diverse anomaly types.}
\response To directly counter the ``alignment with clique injection'' concern, we added a \textbf{non-clique} structural injection (DICE-structural, which rewires edges rather than inserting a dense clique) on a \textbf{fourth dataset (CiteSeer)}, evaluated across \textbf{seven injection ratios (0.5\%--8\%)}. Minomaly detects this different topology reliably (AUC $89$--$97\%$), showing the result is not an artifact of clique-shaped anomalies at one ratio. We also added a \textbf{real-world fraud experiment} (see R3.W4 below) with organic labels. We agree that heterogeneous-graph benchmarks are important and now state them explicitly as future work.

\subsection*{R3.W4 -- Missing dense-subgraph fraud-detection literature}
\comment{Fraudar (KDD 2016), Spade (VLDB 2023), and RUSH (VLDB 2024) detect dense fraudulent blocks without node attributes --- directly matching this setting --- and should be discussed/compared rather than weak baselines like GCNAE.}
\response We thank the reviewer for these precise pointers. We addressed this concern in two ways: positioning these methods in the related work, and adding a real-world fraud experiment that validates Minomaly on organic labels.

\medskip
\textbf{Related work.} We added the dense-subgraph fraud line: \textbf{Fraudar} \citep{hooi2016fraudar} (greedy dense-block peeling without attributes), and its real-time/dynamic successors \textbf{Spade} \citep{jiang2023spade} and \textbf{RUSH} \citep{chen2024rush}. We position Minomaly against them: Fraudar, Spade, and RUSH optimize a global density objective over node blocks, whereas Minomaly measures how rare each node's local neighborhood is in an order-embedding space and returns the specific rare substructure, providing node-level interpretability. Spade and RUSH are also the natural prior art for the dynamic-graph extension that we discuss as future work.

\medskip
\textbf{Real-world fraud experiment.} To our knowledge, no public graph anomaly detection benchmark provides labels that explicitly distinguish structural anomalies from contextual or behavioral ones. To evaluate Minomaly in a setting where labels are organic rather than injected, we applied it to the \textbf{Amazon fraud detection graph} \citep{dou2020enhancing}. This graph models the co-purchase (U-P-U) relation among Amazon reviewers: two users are connected if they review the same product. It contains 11{,}944 nodes, 351{,}216 edges, and 821 labeled fraud nodes (6.9\%).

Since no structural anomaly labels exist, we analyzed the fraud subgraph ourselves. We extracted connected components among fraud nodes and computed each component's internal edge density. Seven components with five or more nodes and density above 0.8 emerge, containing \textbf{136 fraud nodes} (16.6\% of all fraud). The largest is a 79-node near-clique with density 0.95, and four components are perfect 10-node or 5-node cliques. The remaining 685 fraud nodes are scattered throughout the graph with no structural signature.

\medskip
\noindent\textit{Raw detection.}\quad Minomaly runs on the full graph with no knowledge of labels, using the DSAN matcher trained on synthetic graphs. The beam search verifies 6{,}765 anchor nodes as belonging to infrequent neighborhoods, of which 439 are fraud (\textbf{53.5\% recall} against all fraud). The fraud cliques are rare enough to fall within the infrequent band and be discovered by the search. Since this dataset has no structural anomaly annotations, the verified set also includes non-fraudulent rare structures, which is expected in any unsupervised setting.

\medskip
\noindent\textit{Structural interpretation.}\quad Minomaly returns each verified pattern as an explicit subgraph, so the analyst can inspect what kind of rare structure was discovered. Dense cliques are a well-known indicator of coordinated fraud; Minomaly provides the structural information needed to identify them. We define two features for every node in the graph:
\begin{itemize}
  \item \emph{Verified-neighbor count} ($\mathrm{vn\_count}$): how many of the node's graph neighbors are themselves verified anchors.
  \item \emph{Verified-neighbor density} ($\mathrm{vn\_density}$): among those verified-anchor neighbors, the fraction of all possible edges that actually exist.
\end{itemize}
A node embedded in a fraud clique has many verified neighbors that are all connected to each other, producing $\mathrm{vn\_density} \approx 1.0$. A benign node may have a few verified neighbors, but they are scattered and unrelated, producing $\mathrm{vn\_density} \approx 0.0$. Table~\ref{tab:fraud_profile_resp} quantifies this separation.

\begin{table}[h]
\footnotesize
\centering
\caption{Structural profile of true-positive (structural fraud) and false-positive (benign) verified anchors on the Amazon fraud graph. Median values reported.}
\label{tab:fraud_profile_resp}
\begin{tabular}{lcccc}
\toprule
& \textbf{Degree} & \textbf{VN count} & \textbf{VN density} & \textbf{Composite score} \\
\midrule
TP anchors (struct.\ fraud, $n{=}40$)  & 418 & 116 & 1.000 & 4.76 \\
FP anchors (benign, $n{=}2{,}535$)      & 8   & 1   & 0.000 & 0.00 \\
\bottomrule
\end{tabular}
\end{table}

The composite score $\log(\mathrm{vn\_count}+1) \times \mathrm{vn\_density}$ achieves an \textbf{AUC of 94.1\%} against the 136 structural fraud nodes. At the top 1{,}000 ranked nodes, 113 of 136 structural fraud nodes are recovered (83\%), including 97\% of the 79-node clique and 100\% of two 10-node cliques. Only the smallest cliques (5 nodes) fall below the ranking threshold.

\noindent\textit{Implications.}\quad This experiment demonstrates three points:

\begin{enumerate}[leftmargin=1.6em,itemsep=3pt]
  \item \textbf{Structural anomalies exist in real-world fraud.} Among 821 organic fraud labels, 16.6\% form dense cliques that are topologically distinguishable from the rest of the graph. These are not an artifact of synthetic injection.
  \item \textbf{Interpretability enables domain-specific filtering.} Minomaly discovers all rare structures in the graph, including the fraud cliques. Because it returns explicit subgraph patterns rather than opaque scores, the analyst can characterize each pattern and filter for the clique topology that indicates coordinated fraud. This interpretation step achieves AUC 94.1\% against the structural fraud nodes. Methods that produce only a score would require additional post-hoc explanation; Minomaly provides the structure directly.
  \item \textbf{Structural and contextual anomalies are complementary.} The remaining 84\% of fraud nodes are behavioral anomalies with no structural signature. Detecting them requires contextual features such as review text, timestamps, or cross-relation patterns. Integrating such features into the order-embedding framework is a natural extension of Minomaly. A complete fraud detection system should address both types, and Minomaly provides the structural foundation.
\end{enumerate}

%==================================================================
\section{Response to Reviewer 2}
\response We thank the reviewer for the detailed suggestions. We address the points that fall within the paper's stated scope and clarify those that extend beyond it.

\begin{itemize}
  \item \textbf{W1 (node attributes / contextual anomalies) and W2 (dynamic graphs).} Minomaly is, by design and by title, a framework for \emph{structural} anomalies in \emph{static} graphs. Integrating node attributes for contextual anomalies and extending to dynamic graphs are valuable directions that we state explicitly as future work; they define a different problem setting rather than a deficiency of the present contribution. The corresponding details (D1--D4) describe extensions to those settings and are, accordingly, out of the current scope.
  \item \textbf{W3 (adaptive hyperparameters; precision on dense graphs).} Minomaly's hyperparameters are interpretable and controllable, which is a deliberate design choice for an interpretable method; an automatic selection mechanism (D5) is future work. The precision behavior on highly connected graphs (D6) is discussed as a known limitation tied to the greedy search.
  \item \textbf{D7 (non-GNN miners gSpan, GASTON, Motivo).} These are frequent-subgraph \emph{miners}, not node-level anomaly detectors; they are already discussed in the Related Work, and their exponential complexity in pattern size is precisely the limitation that motivates a learned order-embedding alternative. The dense-subgraph fraud detectors added for R3.W4 are the more directly comparable baselines.
  \item \textbf{D8 (more / larger / heterogeneous datasets).} We added CiteSeer (non-clique injection, seven ratios) and the Amazon fraud graph (organic labels); the largest existing dataset (Flickr) already stresses scalability, and ultra-large and heterogeneous graphs are stated as future work.
\end{itemize}

Respectfully, we note that this review is framed largely as proposals for future extensions (``we propose extending\dots'', ``we recommend\dots'') rather than as an assessment of the present contribution, and that most of its points concern problem settings (attributed, dynamic, ultra-large, heterogeneous graphs) outside the paper's stated structural/static scope. We have nonetheless acted on every item that intersects the paper's scope, and we respectfully ask that the decision rest primarily on the substantive technical reviews (Reviewers~1 and~3), whose concerns we have addressed in full.

%==================================================================
\vspace{8pt}\noindent\rule{\textwidth}{0.4pt}

\bibliographystyle{plainnat}
\begin{thebibliography}{4}

\bibitem[Dou et~al.(2020)]{dou2020enhancing}
Y.~Dou, Z.~Liu, L.~Sun, Y.~Deng, H.~Peng, and P.~S. Yu.
\newblock Enhancing graph neural network-based fraud detectors against
  camouflaged fraudsters.
\newblock In \emph{Proceedings of the 29th ACM International Conference on
  Information and Knowledge Management (CIKM)}, pages 315--324, 2020.

\bibitem[Hooi et~al.(2016)]{hooi2016fraudar}
B.~Hooi, H.~A. Song, A.~Beutel, N.~Shah, K.~Shin, and C.~Faloutsos.
\newblock {FRAUDAR}: Bounding graph fraud in the face of camouflage.
\newblock In \emph{Proceedings of the 22nd ACM SIGKDD International Conference
  on Knowledge Discovery and Data Mining}, pages 895--904, 2016.

\bibitem[Jiang et~al.(2023)]{jiang2023spade}
J.~Jiang, J.~Chen, K.~Shin, and C.~Faloutsos.
\newblock {Spade}: Real-time fraud detection on dynamic graphs.
\newblock \emph{Proceedings of the VLDB Endowment}, 16(11):3167--3179, 2023.

\bibitem[Chen et~al.(2024)]{chen2024rush}
J.~Chen, J.~Jiang, K.~Shin, and C.~Faloutsos.
\newblock {RUSH}: Real-time dense subgraph detection on dynamic graphs.
\newblock \emph{Proceedings of the VLDB Endowment}, 17(11):7058--7073, 2024.

\end{thebibliography}

\end{document}
